\section*{Step 273: Chowla Conjecture Classification}

\paragraph{Carrier.}
With distinct shifts $a_1,\dots,a_k$ and standard Mobius exponents
$e_i\in\{1,2\}$, define
\[
  \Xi_{\rm Chowla}(a,e;N)=
  \frac1N\left|\sum_{n\leq N}\prod_{i=1}^k \mu(n+a_i)^{e_i}\right|.
\]
The Chowla conjecture asserts $\Xi_{\rm Chowla}(a,e;N)\to 0$ for all
admissible nontrivial patterns.

\paragraph{Convention note.}
The notation $\epsilon_i=\pm1$ is not literal for the Mobius function, since
$\mu(n)^{-1}$ is undefined when $\mu(n)=0$.  The standard Chowla convention
uses powers $e_i\in\{1,2\}$, with not all exponents even.  For the Liouville
function one can use sign-style powers harmlessly.

\paragraph{Classification.}
Chowla is a pure arithmetic multi-correlation of shifted values of one
multiplicative function.  It fits Type Ia of the Cross-Correlation Extension,
but it is not merely a pairwise coefficient correlation.  Step 273 refines
Type Ia-pure-arithmetic into:
\[
\begin{array}{c|l}
\text{sub-sub-type} & \text{examples}\\\hline
2\text{-correlation} & \text{SOC, full Selberg orthonormality}\\
k\text{-correlation / multi-shift} & \text{Chowla, Elliott}
\end{array}
\]

\paragraph{Known cases.}
Tao's 2016 work proves logarithmically averaged two-point Chowla/Elliott
cases.  Tao--Teravainen prove logarithmically averaged odd-order cases.  The
full ordinary Chowla conjecture remains open.

\paragraph{Verdict.}
\[
  \boxed{\texttt{V\_chowla\_in\_Ia\_pure\_arithmetic\_k\_correlation}}.
\]
No Chowla or RH proof is claimed.
